\section{Enumerations}\label{app:counts}

\subsection{Active-family census at
\texorpdfstring{$(6,3)$ and $(7,3)$}{(6,3) and (7,3)}}\label{app:census}

By \eqref{eq:coverage} the family $A^{+}$ of active $k$-subsets carrying a
positive multiplier covers the index set.  Proposition~\ref{prop:census}
counts the covering families of $3$-subsets up to the symmetric group of
the index set; their distribution by cardinality is recorded here.

\begin{table}[ht]
\small
\centering
\begin{tabular}{@{}l*{10}{r}@{}}
\toprule
$j$   & 2 & 3 & 4 & 5 & 6 & 7 & 8 & 9 & 10 & 11\\
$N_j$ & 1 & 3 & 15 & 37 & 88 & 157 & 247 & 311 & 351 & 312\\
\midrule
$j$   & 12 & 13 & 14 & 15 & 16 & 17 & 18 & 19 & 20 & \\
$N_j$ & 249 & 161 & 94 & 43 & 21 & 7 & 3 & 1 & 1 & \\
\bottomrule
\end{tabular}
\medskip
\caption{The number $N_j$ of families of $j$ triples covering
$\{1,\dots,6\}$, up to relabelling.  The total is $2102$; the
Carath\'eodory bound \eqref{eq:carath} discards the last five columns,
leaving $2069$.}
\label{tab:census}
\end{table}

Setting aside the partition class that Corollary~\ref{cor:e3} leaves
undecided, the columns up to fifteen that bear on
Conjecture~\ref{conj:sixthree} account for $2068$ classes.

At $m=7$ a covering family has at least three members, and the counts are
\[
\begin{array}{c|cccccccc}
\lvert A^{+}\rvert & 3&4&5&6&7&8&9&10\\ \hline
\text{classes} & 3&17&94&415&1599&5308&15397&39081\\[4pt]
\lvert A^{+}\rvert & 11&12&13&14&15&16&17&18\\ \hline
\text{classes} & 87347&172684&303116&473733&660907&824389&920439&920443\\[4pt]
\lvert A^{+}\rvert & 19&20&21&22&23&24&25&26\\ \hline
\text{classes} & 824409&660949&473827&303277&172933&87659&39433&15709\\[4pt]
\lvert A^{+}\rvert & 27&28&29&30&31&32&33&34\\ \hline
\text{classes} & 5557&1760&509&137&38&10&3&1
\end{array}
\]
together with a single class of cardinality $35$.  Enumeration at this
size is not a route; Problem~\ref{prob:seventhree} asks instead for an
argument uniform in the active family.
